%% kapa: exam=1 round=22 year=2011 period=1 form=A subject=civil_law subject_ko=민법 seq=20 number=20
%% source: Q-Net 감정평가사 기출문제 (cst00309) artlSeq=5036790 file=extracted/1008/5036790/2011년 제22회 감정평가사 1차시험 기출문제/1교시.hwp
%% conversion: kordoc 4.24.51 -> tools/gen1.py (auto); review=proofread-sonnet; flags=none
\begin{question}{20}
지상권에 관한 설명으로 옳지 않은 것은? (다툼이 있으면 판례에 의함)
\begin{choices}
  \item 지상권 양도금지특약이 지상권자에게 불리한 경우에는 그 효력이 없다.
  \item 지상권이 존속기간의 만료로 소멸하더라도 지상물이 현존하는 때에는 지상권자는 계약갱신을 청구할 수 있다.
  \item 건물 이외의 공작물의 소유를 목적으로 하는 때에도 30년의 존속기간을 설정할 수 있다.
  \item 지상권이 소멸한 때에 지상권자가 지상물을 수거하고자 하는 경우, 지상권설정자는 상당한 가액을 제공하고 그 지상물의 매수를 청구할 수 있다.
  \item 지상권의 지료가 1년 연체된 상태에서 토지가 제3자에게 양도되고 다시 그 지료가 1년 6개월 연체된 경우, 토지의 새로운 소유자는 지상권의 소멸을 청구할 수 있다.
\end{choices}
\end{question}
\answer{5}
